% ==============================================================================
% PDF-1 / Section 01 — Scope & Reproducibility (Skeleton)
% No new primitives. Non-prescriptive. Container only.
% ==============================================================================

\section{Scope and Reproducibility}
\label{sec:scope-repro}

\subsection{Document class and intent}
This document is an \emph{executable demonstration} of diagnostic \emph{phase space}
construction and \emph{STOP regions} for the Enterprise Continuum \(K_{EA}\).
It is strictly \emph{non-prescriptive}: it does not recommend interventions, policies,
or managerial actions.

The present PDF-1 is a \emph{container skeleton}. It provides a stable LaTeX structure
for later computational demos (in subsequent runs), while preserving the immutability of:
(i) KEA (PDF-0) and (ii) \texttt{ETS.K\_EA.v1.1}.
No new theoretical primitives are introduced here, and no semantic drift is permitted.

\subsection{Immutability and scope guardrails}
\begin{itemize}
  \item \textbf{Immutable inputs:} KEA (PDF-0) and \texttt{ETS.K\_EA.v1.1} are treated as authoritative.
  \item \textbf{No new primitives:} only re-use of existing definitions and terms.
  \item \textbf{No engine in RUN-EDHQ-01:} no computation, no implementation details, no algorithms beyond placeholders.
  \item \textbf{No pipeline edits:} the shared preamble and build tooling remain untouched.
\end{itemize}

\subsection{Reproducible build path}
The build must succeed when executed from the repository root using PowerShell and \texttt{latexmk}.
The canonical local build commands are:

\begin{verbatim}
Set-Location C:\Users\Megaport\work\ea-oc-extension
Push-Location .\paper\pdf1_executing_kea
latexmk -pdf -interaction=nonstopmode -file-line-error -halt-on-error `
  .\main.tex
Pop-Location
\end{verbatim}

Expected output:
\begin{itemize}
  \item PDF artifact: \path{paper/pdf1_executing_kea/main.pdf}
  \item Exit status: \texttt{0} (success)
\end{itemize}

\subsection{What this PDF-1 will contain (structural)}
Subsequent sections provide only the structural slots for:
(i) phase space construction, (ii) synthetic regimes, (iii) executing diagnostics,
(iv) STOP regions, (v) invariance stress tests, and (vi) explicit non-claims and limits.
Any future computational content must remain diagnostic-only and must not modify the KEA/ETS core.
